% !TeX root = ../../../金融经济学_王江_讲义及习题.tex
% 左图在权重空间中取中点，右图比较实际风险与端点标准差的算术平均。
\begin{tikzpicture}[
 >=Stealth,font=\small,
 frontier/.style={blue!70!black,thick},
 mix/.style={teal!65!black},
 guide/.style={gray,densely dotted}]
 \begin{scope}[x=1.45cm,y=1cm]
  \coordinate (A) at (.25,.73);
  \coordinate (B) at (2.65,2.53);
  \coordinate (C) at (1.45,1.63);
  \draw[frontier,<->] (-.07,.49)--(2.97,2.77);
  \fill[blue!70!black] (A) circle (1.8pt) node[above left] {$z_A$};
  \fill[blue!70!black] (B) circle (1.8pt) node[above left] {$z_B$};
  \fill[mix] (C) circle (2pt) node[below right=2pt] {$z_C=\dfrac{z_A+z_B}{2}$};
  \node[blue!70!black] at (1.12,2.77) {前沿直线};
  \node at (1.55,-.40) {(a) 权重空间};
 \end{scope}
 \begin{scope}[xshift=6.25cm,x=1.60cm,y=.90cm]
  \draw[->] (0,0)--(2.70,0) node[right] {$\sigma$};
  \draw[->] (0,0)--(0,3.58) node[above] {$\bar r$};
  % 示意双曲线 x^2-(y-1.45)^2=1；三个点的纵坐标依次为 .55、1.75、2.95。
  \coordinate (A) at ({sqrt(1+.9*.9)},.55);
  \coordinate (B) at ({sqrt(1+1.5*1.5)},2.95);
  \coordinate (C) at ({sqrt(1+.3*.3)},1.75);
  \coordinate (chordmid) at ({(sqrt(1+.9*.9)+sqrt(1+1.5*1.5))/2},1.75);
  \draw[guide] (0,.55) node[left,black] {$\bar r_A$}--(A);
  \draw[guide] (0,2.95) node[left,black] {$\bar r_B$}--(B);
  \draw[guide] (0,1.75) node[left,black] {$\bar r_C$}--(chordmid);
  \draw[gray,dashed] (A)--(B);
  \draw[frontier,domain=-1.20:1.80,samples=81,variable=\t]
   plot ({sqrt(1+\t*\t)},{1.45+\t});
  \fill[blue!70!black] (A) circle (1.8pt) node[right=2pt] {$A$};
  \fill[blue!70!black] (B) circle (1.8pt) node[right=2pt] {$B$};
  \fill[mix] (C) circle (2pt) node[above left=1pt] {$C$};
  \draw[gray,fill=white] (chordmid) circle (1.8pt)
   node[right=2pt,align=left] {标准差\\取平均};
  \node at (1.38,-.44) {(b) 均值--标准差平面};
 \end{scope}
\end{tikzpicture}
